% Methods section for Pyro-DCM paper submission.
% This is a fragment -- include via % Methods section for Pyro-DCM paper submission.
% This is a fragment -- include via % Methods section for Pyro-DCM paper submission.
% This is a fragment -- include via % Methods section for Pyro-DCM paper submission.
% This is a fragment -- include via \input{methods.tex} in your paper template.
% Requires: amsmath, amssymb, booktabs packages.
% Bibliography: references.bib (see companion file).

\section{Methods}

\subsection{Generative Model}

\subsubsection{Neural State Equation}

The neural dynamics of the DCM are described by the bilinear state
equation \citep{friston2003}:
\begin{equation}
  \frac{dx}{dt} = Ax + Cu(t)
  \label{eq:neural-state}
\end{equation}
where $x \in \mathbb{R}^N$ is the neural activity vector for $N$ brain
regions, $A \in \mathbb{R}^{N \times N}$ is the effective connectivity
matrix encoding directed causal influences between regions, $C \in
\mathbb{R}^{N \times M}$ maps $M$ experimental inputs $u(t)$ to regions,
and $t$ denotes time. The diagonal elements $a_{ii}$ represent
self-inhibition, and off-diagonal elements $a_{ij}$ represent the
influence of region $j$ on region $i$.

The $A$ matrix is parameterized following SPM12 conventions to guarantee
stability:
\begin{equation}
  a_{ii} = -\exp(A^\text{free}_{ii}) / 2, \qquad
  a_{ij} = A^\text{free}_{ij} \quad (i \neq j)
  \label{eq:a-param}
\end{equation}
ensuring negative self-connections (default $-0.5$~Hz when the free
parameter is zero).

\subsubsection{Hemodynamic Model}

The Balloon-Windkessel model \citep{stephan2007} translates neural
activity into hemodynamic responses through four coupled ordinary
differential equations per region, implemented in log-space following
SPM12 (\texttt{spm\_fx\_fmri.m}):
\begin{align}
  \frac{ds}{dt} &= x - \kappa s - \gamma (f - 1)
    \label{eq:vasodilatory} \\
  \frac{d \ln f}{dt} &= \frac{s}{f}
    \label{eq:flow} \\
  \frac{d \ln v}{dt} &= \frac{f - v^{1/\alpha}}{\tau v}
    \label{eq:volume} \\
  \frac{d \ln q}{dt} &= \frac{f\, E(f, E_0) / E_0 - v^{1/\alpha} q / v}{\tau q}
    \label{eq:deoxyhemoglobin}
\end{align}
where $E(f, E_0) = 1 - (1 - E_0)^{1/f}$ is the oxygen extraction
fraction.

\begin{table}[h]
  \centering
  \caption{Hemodynamic parameter defaults (SPM12 code).}
  \label{tab:hemo-params}
  \begin{tabular}{llccl}
    \toprule
    Parameter & Symbol & Default & Unit & Description \\
    \midrule
    Signal decay       & $\kappa$ & 0.64 & s$^{-1}$ & Rate of vasodilatory signal decay \\
    Autoregulation     & $\gamma$ & 0.32 & s$^{-1}$ & Flow-dependent elimination \\
    Transit time       & $\tau$   & 2.00 & s        & Hemodynamic transit time \\
    Grubb's exponent   & $\alpha$ & 0.32 & --       & Vessel stiffness \\
    O$_2$ extraction   & $E_0$    & 0.40 & --       & Resting oxygen extraction fraction \\
    \bottomrule
  \end{tabular}
\end{table}

\subsubsection{BOLD Signal}

The BOLD percent signal change is computed from hemodynamic states via
the simplified Buxton observation equation \citep{stephan2007}:
\begin{equation}
  y = V_0 \left[ k_1 (1 - q) + k_2 \left(1 - \frac{q}{v}\right) + k_3 (1 - v) \right]
  \label{eq:bold}
\end{equation}
where $k_1 = 7 E_0$, $k_2 = 2$, $k_3 = 2 E_0 - 0.2$, and
$V_0 = 0.02$ is the resting venous blood volume fraction.

\subsubsection{Spectral DCM}

For resting-state fMRI, spectral DCM operates in the frequency domain
using cross-spectral density (CSD) as the data representation
\citep{friston2014}.

The transfer function is computed via eigendecomposition:
\begin{equation}
  H(\omega) = C_\text{out} (i\omega I - A)^{-1} C_\text{in}
    = \sum_k \frac{\mathbf{g}_k \mathbf{u}_k^T}{i 2\pi\omega - \lambda_k}
  \label{eq:transfer}
\end{equation}
where $\lambda_k$ are the eigenvalues of $A$ (stabilized by clamping
$\text{Re}(\lambda_k) \leq -1/32$) and $C_\text{in} = C_\text{out} = I_N$
in the standard spDCM convention.

The predicted CSD is:
\begin{equation}
  S(\omega) = H(\omega)\, G_u(\omega)\, H(\omega)^\dagger + G_n(\omega)
  \label{eq:predicted-csd}
\end{equation}
where the neuronal noise follows a $1/f^\alpha$ power law:
\begin{equation}
  G_{u,i}(\omega) = C \cdot \exp(a_{0,i}) \cdot \omega^{-\exp(a_{1,i})} \cdot 4
  \label{eq:neuronal-noise}
\end{equation}
with $C = 1/256$ (SPM scaling constant), and the observation noise
combines global and regional components:
\begin{equation}
  G_n(\omega) = G_\text{global}(\omega) \cdot \mathbf{1}\mathbf{1}^T
    + \text{diag}(G_\text{regional}(\omega))
  \label{eq:obs-noise}
\end{equation}

\subsubsection{Regression DCM}

Regression DCM \citep{frassle2017} reformulates DCM as a per-region
linear regression in the frequency domain:
\begin{equation}
  y_j = X_j \theta_j + \varepsilon_j
  \label{eq:rdcm-regression}
\end{equation}
where $y_j$ is the DFT of region $j$'s BOLD signal (real and imaginary
parts stacked), $X_j$ is the frequency-domain design matrix containing
HRF-convolved regressors from other regions and external inputs, and
$\theta_j$ contains the connectivity and input parameters for region $j$.

\subsection{Inference}

\subsubsection{Stochastic Variational Inference}

We replace the Variational Laplace algorithm used in SPM12
\citep{friston2007vfe} with stochastic variational inference (SVI;
\citealp{blei2017}) implemented in Pyro \citep{bingham2019}. The ELBO
objective is:
\begin{equation}
  \mathcal{L}(\phi)
    = \mathbb{E}_{q_\phi(\theta)}
      \left[ \log p(y, \theta \mid m) \right]
    - \mathbb{E}_{q_\phi(\theta)}
      \left[ \log q_\phi(\theta) \right]
  \label{eq:elbo}
\end{equation}
where the variational family is a mean-field Gaussian:
\begin{equation}
  q_\phi(\theta) = \prod_i \mathcal{N}(\theta_i;\, \mu_i, \sigma_i^2)
  \label{eq:mean-field}
\end{equation}
Optimization uses ClippedAdam with initial learning rate 0.01, gradient
clipping norm 10.0, and exponential learning rate decay to 1\% over the
training run. ODE-based models use fixed-step rk4 integration with
$\Delta t = 0.5$~s during SVI.

Priors: $A_\text{free} \sim \mathcal{N}(0, 1/64)$ (SPM12 convention),
$C \sim \mathcal{N}(0, 1)$.

\subsubsection{Amortized Inference}

For rapid posterior estimation across subjects, we provide amortized
inference using Neural Spline Flows \citep{papamakarios2021} conditioned
on summary statistics \citep{cranmer2020}:
\begin{equation}
  q_\phi(\theta \mid \mathbf{x})
    = f_\phi\!\left(\mathbf{z};\, \text{summary}(\mathbf{x})\right)
  \label{eq:amortized}
\end{equation}
Summary networks compress BOLD time series (1D-CNN) or CSD matrices
(MLP) into fixed-dimensional embeddings. Parameters are packed into a
standardized vector on $[-5, 5]$ matching the NSF spline domain.

\subsubsection{Analytic VB for rDCM}

Regression DCM admits closed-form posterior computation
\citep{frassle2017}:
\begin{align}
  \Sigma_j &= \left(\Lambda_0 + \langle\lambda_j\rangle X_j^H X_j\right)^{-1}
    \label{eq:rdcm-cov} \\
  \mu_j &= \Sigma_j
    \left(\Lambda_0 m_0 + \langle\lambda_j\rangle X_j^H y_j\right)
    \label{eq:rdcm-mean}
\end{align}
with noise precision updates $a_N = a_0 + N_y/2$ and
$b_N = b_0 + \frac{1}{2}(y^H y - \mu^T \Sigma^{-1} \mu + m_0^T \Lambda_0 m_0)$.

The free energy is computed per region as the sum of five terms
(log-likelihood, prior complexity, posterior entropy, noise prior, noise
posterior). The sparse variant \citep{frassle2018} extends this with ARD
binary indicators $z_{j,k} \sim \text{Bernoulli}(0.5)$.

\subsection{Bayesian Model Comparison}

For task and spectral DCM, the final ELBO serves as a lower bound on the
log model evidence \citep{friston2007vfe}. The model with the higher
ELBO is preferred. For rDCM, the analytic free energy $F = \sum_r F_r$
provides an exact bound within the variational approximation
\citep{frassle2017}. A difference $\Delta > 3$~nats (Bayes factor $> 20$)
is considered strong evidence.

\subsection{Implementation}

Pyro-DCM is implemented in Python 3.10+ using PyTorch 2.x for tensor
computation, Pyro 1.9+ for probabilistic programming, torchdiffeq for
ODE integration, and Zuko for normalizing flows. Numerical stability is
ensured via log-space hemodynamic states, eigenvalue stabilization
($\text{Re}(\lambda) \leq -1/32$), gradient clipping (norm 10.0),
NaN detection with early termination, and \texttt{torch.linalg.solve}
for matrix operations.

ODE integration uses adaptive dopri5 for simulation ($\Delta t = 0.01$~s)
and fixed-step rk4 for SVI ($\Delta t = 0.5$~s) to ensure predictable
runtime during optimization.

\subsection{Benchmarks}

Parameter recovery benchmarks follow a simulate--infer--compare protocol
with 3-region networks. Metrics include RMSE$(A)$, 90\% coverage,
Pearson correlation, ELBO/free energy, and wall time per subject.
Amortization-specific metrics include the RMSE ratio, amortization gap
($\text{ELBO}_\text{SVI} - \text{ELBO}_\text{amortized}$), and speed
ratio. See the benchmark report for detailed results.
 in your paper template.
% Requires: amsmath, amssymb, booktabs packages.
% Bibliography: references.bib (see companion file).

\section{Methods}

\subsection{Generative Model}

\subsubsection{Neural State Equation}

The neural dynamics of the DCM are described by the bilinear state
equation \citep{friston2003}:
\begin{equation}
  \frac{dx}{dt} = Ax + Cu(t)
  \label{eq:neural-state}
\end{equation}
where $x \in \mathbb{R}^N$ is the neural activity vector for $N$ brain
regions, $A \in \mathbb{R}^{N \times N}$ is the effective connectivity
matrix encoding directed causal influences between regions, $C \in
\mathbb{R}^{N \times M}$ maps $M$ experimental inputs $u(t)$ to regions,
and $t$ denotes time. The diagonal elements $a_{ii}$ represent
self-inhibition, and off-diagonal elements $a_{ij}$ represent the
influence of region $j$ on region $i$.

The $A$ matrix is parameterized following SPM12 conventions to guarantee
stability:
\begin{equation}
  a_{ii} = -\exp(A^\text{free}_{ii}) / 2, \qquad
  a_{ij} = A^\text{free}_{ij} \quad (i \neq j)
  \label{eq:a-param}
\end{equation}
ensuring negative self-connections (default $-0.5$~Hz when the free
parameter is zero).

\subsubsection{Hemodynamic Model}

The Balloon-Windkessel model \citep{stephan2007} translates neural
activity into hemodynamic responses through four coupled ordinary
differential equations per region, implemented in log-space following
SPM12 (\texttt{spm\_fx\_fmri.m}):
\begin{align}
  \frac{ds}{dt} &= x - \kappa s - \gamma (f - 1)
    \label{eq:vasodilatory} \\
  \frac{d \ln f}{dt} &= \frac{s}{f}
    \label{eq:flow} \\
  \frac{d \ln v}{dt} &= \frac{f - v^{1/\alpha}}{\tau v}
    \label{eq:volume} \\
  \frac{d \ln q}{dt} &= \frac{f\, E(f, E_0) / E_0 - v^{1/\alpha} q / v}{\tau q}
    \label{eq:deoxyhemoglobin}
\end{align}
where $E(f, E_0) = 1 - (1 - E_0)^{1/f}$ is the oxygen extraction
fraction.

\begin{table}[h]
  \centering
  \caption{Hemodynamic parameter defaults (SPM12 code).}
  \label{tab:hemo-params}
  \begin{tabular}{llccl}
    \toprule
    Parameter & Symbol & Default & Unit & Description \\
    \midrule
    Signal decay       & $\kappa$ & 0.64 & s$^{-1}$ & Rate of vasodilatory signal decay \\
    Autoregulation     & $\gamma$ & 0.32 & s$^{-1}$ & Flow-dependent elimination \\
    Transit time       & $\tau$   & 2.00 & s        & Hemodynamic transit time \\
    Grubb's exponent   & $\alpha$ & 0.32 & --       & Vessel stiffness \\
    O$_2$ extraction   & $E_0$    & 0.40 & --       & Resting oxygen extraction fraction \\
    \bottomrule
  \end{tabular}
\end{table}

\subsubsection{BOLD Signal}

The BOLD percent signal change is computed from hemodynamic states via
the simplified Buxton observation equation \citep{stephan2007}:
\begin{equation}
  y = V_0 \left[ k_1 (1 - q) + k_2 \left(1 - \frac{q}{v}\right) + k_3 (1 - v) \right]
  \label{eq:bold}
\end{equation}
where $k_1 = 7 E_0$, $k_2 = 2$, $k_3 = 2 E_0 - 0.2$, and
$V_0 = 0.02$ is the resting venous blood volume fraction.

\subsubsection{Spectral DCM}

For resting-state fMRI, spectral DCM operates in the frequency domain
using cross-spectral density (CSD) as the data representation
\citep{friston2014}.

The transfer function is computed via eigendecomposition:
\begin{equation}
  H(\omega) = C_\text{out} (i\omega I - A)^{-1} C_\text{in}
    = \sum_k \frac{\mathbf{g}_k \mathbf{u}_k^T}{i 2\pi\omega - \lambda_k}
  \label{eq:transfer}
\end{equation}
where $\lambda_k$ are the eigenvalues of $A$ (stabilized by clamping
$\text{Re}(\lambda_k) \leq -1/32$) and $C_\text{in} = C_\text{out} = I_N$
in the standard spDCM convention.

The predicted CSD is:
\begin{equation}
  S(\omega) = H(\omega)\, G_u(\omega)\, H(\omega)^\dagger + G_n(\omega)
  \label{eq:predicted-csd}
\end{equation}
where the neuronal noise follows a $1/f^\alpha$ power law:
\begin{equation}
  G_{u,i}(\omega) = C \cdot \exp(a_{0,i}) \cdot \omega^{-\exp(a_{1,i})} \cdot 4
  \label{eq:neuronal-noise}
\end{equation}
with $C = 1/256$ (SPM scaling constant), and the observation noise
combines global and regional components:
\begin{equation}
  G_n(\omega) = G_\text{global}(\omega) \cdot \mathbf{1}\mathbf{1}^T
    + \text{diag}(G_\text{regional}(\omega))
  \label{eq:obs-noise}
\end{equation}

\subsubsection{Regression DCM}

Regression DCM \citep{frassle2017} reformulates DCM as a per-region
linear regression in the frequency domain:
\begin{equation}
  y_j = X_j \theta_j + \varepsilon_j
  \label{eq:rdcm-regression}
\end{equation}
where $y_j$ is the DFT of region $j$'s BOLD signal (real and imaginary
parts stacked), $X_j$ is the frequency-domain design matrix containing
HRF-convolved regressors from other regions and external inputs, and
$\theta_j$ contains the connectivity and input parameters for region $j$.

\subsection{Inference}

\subsubsection{Stochastic Variational Inference}

We replace the Variational Laplace algorithm used in SPM12
\citep{friston2007vfe} with stochastic variational inference (SVI;
\citealp{blei2017}) implemented in Pyro \citep{bingham2019}. The ELBO
objective is:
\begin{equation}
  \mathcal{L}(\phi)
    = \mathbb{E}_{q_\phi(\theta)}
      \left[ \log p(y, \theta \mid m) \right]
    - \mathbb{E}_{q_\phi(\theta)}
      \left[ \log q_\phi(\theta) \right]
  \label{eq:elbo}
\end{equation}
where the variational family is a mean-field Gaussian:
\begin{equation}
  q_\phi(\theta) = \prod_i \mathcal{N}(\theta_i;\, \mu_i, \sigma_i^2)
  \label{eq:mean-field}
\end{equation}
Optimization uses ClippedAdam with initial learning rate 0.01, gradient
clipping norm 10.0, and exponential learning rate decay to 1\% over the
training run. ODE-based models use fixed-step rk4 integration with
$\Delta t = 0.5$~s during SVI.

Priors: $A_\text{free} \sim \mathcal{N}(0, 1/64)$ (SPM12 convention),
$C \sim \mathcal{N}(0, 1)$.

\subsubsection{Amortized Inference}

For rapid posterior estimation across subjects, we provide amortized
inference using Neural Spline Flows \citep{papamakarios2021} conditioned
on summary statistics \citep{cranmer2020}:
\begin{equation}
  q_\phi(\theta \mid \mathbf{x})
    = f_\phi\!\left(\mathbf{z};\, \text{summary}(\mathbf{x})\right)
  \label{eq:amortized}
\end{equation}
Summary networks compress BOLD time series (1D-CNN) or CSD matrices
(MLP) into fixed-dimensional embeddings. Parameters are packed into a
standardized vector on $[-5, 5]$ matching the NSF spline domain.

\subsubsection{Analytic VB for rDCM}

Regression DCM admits closed-form posterior computation
\citep{frassle2017}:
\begin{align}
  \Sigma_j &= \left(\Lambda_0 + \langle\lambda_j\rangle X_j^H X_j\right)^{-1}
    \label{eq:rdcm-cov} \\
  \mu_j &= \Sigma_j
    \left(\Lambda_0 m_0 + \langle\lambda_j\rangle X_j^H y_j\right)
    \label{eq:rdcm-mean}
\end{align}
with noise precision updates $a_N = a_0 + N_y/2$ and
$b_N = b_0 + \frac{1}{2}(y^H y - \mu^T \Sigma^{-1} \mu + m_0^T \Lambda_0 m_0)$.

The free energy is computed per region as the sum of five terms
(log-likelihood, prior complexity, posterior entropy, noise prior, noise
posterior). The sparse variant \citep{frassle2018} extends this with ARD
binary indicators $z_{j,k} \sim \text{Bernoulli}(0.5)$.

\subsection{Bayesian Model Comparison}

For task and spectral DCM, the final ELBO serves as a lower bound on the
log model evidence \citep{friston2007vfe}. The model with the higher
ELBO is preferred. For rDCM, the analytic free energy $F = \sum_r F_r$
provides an exact bound within the variational approximation
\citep{frassle2017}. A difference $\Delta > 3$~nats (Bayes factor $> 20$)
is considered strong evidence.

\subsection{Implementation}

Pyro-DCM is implemented in Python 3.10+ using PyTorch 2.x for tensor
computation, Pyro 1.9+ for probabilistic programming, torchdiffeq for
ODE integration, and Zuko for normalizing flows. Numerical stability is
ensured via log-space hemodynamic states, eigenvalue stabilization
($\text{Re}(\lambda) \leq -1/32$), gradient clipping (norm 10.0),
NaN detection with early termination, and \texttt{torch.linalg.solve}
for matrix operations.

ODE integration uses adaptive dopri5 for simulation ($\Delta t = 0.01$~s)
and fixed-step rk4 for SVI ($\Delta t = 0.5$~s) to ensure predictable
runtime during optimization.

\subsection{Benchmarks}

Parameter recovery benchmarks follow a simulate--infer--compare protocol
with 3-region networks. Metrics include RMSE$(A)$, 90\% coverage,
Pearson correlation, ELBO/free energy, and wall time per subject.
Amortization-specific metrics include the RMSE ratio, amortization gap
($\text{ELBO}_\text{SVI} - \text{ELBO}_\text{amortized}$), and speed
ratio. See the benchmark report for detailed results.
 in your paper template.
% Requires: amsmath, amssymb, booktabs packages.
% Bibliography: references.bib (see companion file).

\section{Methods}

\subsection{Generative Model}

\subsubsection{Neural State Equation}

The neural dynamics of the DCM are described by the bilinear state
equation \citep{friston2003}:
\begin{equation}
  \frac{dx}{dt} = Ax + Cu(t)
  \label{eq:neural-state}
\end{equation}
where $x \in \mathbb{R}^N$ is the neural activity vector for $N$ brain
regions, $A \in \mathbb{R}^{N \times N}$ is the effective connectivity
matrix encoding directed causal influences between regions, $C \in
\mathbb{R}^{N \times M}$ maps $M$ experimental inputs $u(t)$ to regions,
and $t$ denotes time. The diagonal elements $a_{ii}$ represent
self-inhibition, and off-diagonal elements $a_{ij}$ represent the
influence of region $j$ on region $i$.

The $A$ matrix is parameterized following SPM12 conventions to guarantee
stability:
\begin{equation}
  a_{ii} = -\exp(A^\text{free}_{ii}) / 2, \qquad
  a_{ij} = A^\text{free}_{ij} \quad (i \neq j)
  \label{eq:a-param}
\end{equation}
ensuring negative self-connections (default $-0.5$~Hz when the free
parameter is zero).

\subsubsection{Hemodynamic Model}

The Balloon-Windkessel model \citep{stephan2007} translates neural
activity into hemodynamic responses through four coupled ordinary
differential equations per region, implemented in log-space following
SPM12 (\texttt{spm\_fx\_fmri.m}):
\begin{align}
  \frac{ds}{dt} &= x - \kappa s - \gamma (f - 1)
    \label{eq:vasodilatory} \\
  \frac{d \ln f}{dt} &= \frac{s}{f}
    \label{eq:flow} \\
  \frac{d \ln v}{dt} &= \frac{f - v^{1/\alpha}}{\tau v}
    \label{eq:volume} \\
  \frac{d \ln q}{dt} &= \frac{f\, E(f, E_0) / E_0 - v^{1/\alpha} q / v}{\tau q}
    \label{eq:deoxyhemoglobin}
\end{align}
where $E(f, E_0) = 1 - (1 - E_0)^{1/f}$ is the oxygen extraction
fraction.

\begin{table}[h]
  \centering
  \caption{Hemodynamic parameter defaults (SPM12 code).}
  \label{tab:hemo-params}
  \begin{tabular}{llccl}
    \toprule
    Parameter & Symbol & Default & Unit & Description \\
    \midrule
    Signal decay       & $\kappa$ & 0.64 & s$^{-1}$ & Rate of vasodilatory signal decay \\
    Autoregulation     & $\gamma$ & 0.32 & s$^{-1}$ & Flow-dependent elimination \\
    Transit time       & $\tau$   & 2.00 & s        & Hemodynamic transit time \\
    Grubb's exponent   & $\alpha$ & 0.32 & --       & Vessel stiffness \\
    O$_2$ extraction   & $E_0$    & 0.40 & --       & Resting oxygen extraction fraction \\
    \bottomrule
  \end{tabular}
\end{table}

\subsubsection{BOLD Signal}

The BOLD percent signal change is computed from hemodynamic states via
the simplified Buxton observation equation \citep{stephan2007}:
\begin{equation}
  y = V_0 \left[ k_1 (1 - q) + k_2 \left(1 - \frac{q}{v}\right) + k_3 (1 - v) \right]
  \label{eq:bold}
\end{equation}
where $k_1 = 7 E_0$, $k_2 = 2$, $k_3 = 2 E_0 - 0.2$, and
$V_0 = 0.02$ is the resting venous blood volume fraction.

\subsubsection{Spectral DCM}

For resting-state fMRI, spectral DCM operates in the frequency domain
using cross-spectral density (CSD) as the data representation
\citep{friston2014}.

The transfer function is computed via eigendecomposition:
\begin{equation}
  H(\omega) = C_\text{out} (i\omega I - A)^{-1} C_\text{in}
    = \sum_k \frac{\mathbf{g}_k \mathbf{u}_k^T}{i 2\pi\omega - \lambda_k}
  \label{eq:transfer}
\end{equation}
where $\lambda_k$ are the eigenvalues of $A$ (stabilized by clamping
$\text{Re}(\lambda_k) \leq -1/32$) and $C_\text{in} = C_\text{out} = I_N$
in the standard spDCM convention.

The predicted CSD is:
\begin{equation}
  S(\omega) = H(\omega)\, G_u(\omega)\, H(\omega)^\dagger + G_n(\omega)
  \label{eq:predicted-csd}
\end{equation}
where the neuronal noise follows a $1/f^\alpha$ power law:
\begin{equation}
  G_{u,i}(\omega) = C \cdot \exp(a_{0,i}) \cdot \omega^{-\exp(a_{1,i})} \cdot 4
  \label{eq:neuronal-noise}
\end{equation}
with $C = 1/256$ (SPM scaling constant), and the observation noise
combines global and regional components:
\begin{equation}
  G_n(\omega) = G_\text{global}(\omega) \cdot \mathbf{1}\mathbf{1}^T
    + \text{diag}(G_\text{regional}(\omega))
  \label{eq:obs-noise}
\end{equation}

\subsubsection{Regression DCM}

Regression DCM \citep{frassle2017} reformulates DCM as a per-region
linear regression in the frequency domain:
\begin{equation}
  y_j = X_j \theta_j + \varepsilon_j
  \label{eq:rdcm-regression}
\end{equation}
where $y_j$ is the DFT of region $j$'s BOLD signal (real and imaginary
parts stacked), $X_j$ is the frequency-domain design matrix containing
HRF-convolved regressors from other regions and external inputs, and
$\theta_j$ contains the connectivity and input parameters for region $j$.

\subsection{Inference}

\subsubsection{Stochastic Variational Inference}

We replace the Variational Laplace algorithm used in SPM12
\citep{friston2007vfe} with stochastic variational inference (SVI;
\citealp{blei2017}) implemented in Pyro \citep{bingham2019}. The ELBO
objective is:
\begin{equation}
  \mathcal{L}(\phi)
    = \mathbb{E}_{q_\phi(\theta)}
      \left[ \log p(y, \theta \mid m) \right]
    - \mathbb{E}_{q_\phi(\theta)}
      \left[ \log q_\phi(\theta) \right]
  \label{eq:elbo}
\end{equation}
where the variational family is a mean-field Gaussian:
\begin{equation}
  q_\phi(\theta) = \prod_i \mathcal{N}(\theta_i;\, \mu_i, \sigma_i^2)
  \label{eq:mean-field}
\end{equation}
Optimization uses ClippedAdam with initial learning rate 0.01, gradient
clipping norm 10.0, and exponential learning rate decay to 1\% over the
training run. ODE-based models use fixed-step rk4 integration with
$\Delta t = 0.5$~s during SVI.

Priors: $A_\text{free} \sim \mathcal{N}(0, 1/64)$ (SPM12 convention),
$C \sim \mathcal{N}(0, 1)$.

\subsubsection{Amortized Inference}

For rapid posterior estimation across subjects, we provide amortized
inference using Neural Spline Flows \citep{papamakarios2021} conditioned
on summary statistics \citep{cranmer2020}:
\begin{equation}
  q_\phi(\theta \mid \mathbf{x})
    = f_\phi\!\left(\mathbf{z};\, \text{summary}(\mathbf{x})\right)
  \label{eq:amortized}
\end{equation}
Summary networks compress BOLD time series (1D-CNN) or CSD matrices
(MLP) into fixed-dimensional embeddings. Parameters are packed into a
standardized vector on $[-5, 5]$ matching the NSF spline domain.

\subsubsection{Analytic VB for rDCM}

Regression DCM admits closed-form posterior computation
\citep{frassle2017}:
\begin{align}
  \Sigma_j &= \left(\Lambda_0 + \langle\lambda_j\rangle X_j^H X_j\right)^{-1}
    \label{eq:rdcm-cov} \\
  \mu_j &= \Sigma_j
    \left(\Lambda_0 m_0 + \langle\lambda_j\rangle X_j^H y_j\right)
    \label{eq:rdcm-mean}
\end{align}
with noise precision updates $a_N = a_0 + N_y/2$ and
$b_N = b_0 + \frac{1}{2}(y^H y - \mu^T \Sigma^{-1} \mu + m_0^T \Lambda_0 m_0)$.

The free energy is computed per region as the sum of five terms
(log-likelihood, prior complexity, posterior entropy, noise prior, noise
posterior). The sparse variant \citep{frassle2018} extends this with ARD
binary indicators $z_{j,k} \sim \text{Bernoulli}(0.5)$.

\subsection{Bayesian Model Comparison}

For task and spectral DCM, the final ELBO serves as a lower bound on the
log model evidence \citep{friston2007vfe}. The model with the higher
ELBO is preferred. For rDCM, the analytic free energy $F = \sum_r F_r$
provides an exact bound within the variational approximation
\citep{frassle2017}. A difference $\Delta > 3$~nats (Bayes factor $> 20$)
is considered strong evidence.

\subsection{Implementation}

Pyro-DCM is implemented in Python 3.10+ using PyTorch 2.x for tensor
computation, Pyro 1.9+ for probabilistic programming, torchdiffeq for
ODE integration, and Zuko for normalizing flows. Numerical stability is
ensured via log-space hemodynamic states, eigenvalue stabilization
($\text{Re}(\lambda) \leq -1/32$), gradient clipping (norm 10.0),
NaN detection with early termination, and \texttt{torch.linalg.solve}
for matrix operations.

ODE integration uses adaptive dopri5 for simulation ($\Delta t = 0.01$~s)
and fixed-step rk4 for SVI ($\Delta t = 0.5$~s) to ensure predictable
runtime during optimization.

\subsection{Benchmarks}

Parameter recovery benchmarks follow a simulate--infer--compare protocol
with 3-region networks. Metrics include RMSE$(A)$, 90\% coverage,
Pearson correlation, ELBO/free energy, and wall time per subject.
Amortization-specific metrics include the RMSE ratio, amortization gap
($\text{ELBO}_\text{SVI} - \text{ELBO}_\text{amortized}$), and speed
ratio. See the benchmark report for detailed results.
 in your paper template.
% Requires: amsmath, amssymb, booktabs packages.
% Bibliography: references.bib (see companion file).

\section{Methods}

\subsection{Generative Model}

\subsubsection{Neural State Equation}

The neural dynamics of the DCM are described by the bilinear state
equation \citep{friston2003}:
\begin{equation}
  \frac{dx}{dt} = Ax + Cu(t)
  \label{eq:neural-state}
\end{equation}
where $x \in \mathbb{R}^N$ is the neural activity vector for $N$ brain
regions, $A \in \mathbb{R}^{N \times N}$ is the effective connectivity
matrix encoding directed causal influences between regions, $C \in
\mathbb{R}^{N \times M}$ maps $M$ experimental inputs $u(t)$ to regions,
and $t$ denotes time. The diagonal elements $a_{ii}$ represent
self-inhibition, and off-diagonal elements $a_{ij}$ represent the
influence of region $j$ on region $i$.

The $A$ matrix is parameterized following SPM12 conventions to guarantee
stability:
\begin{equation}
  a_{ii} = -\exp(A^\text{free}_{ii}) / 2, \qquad
  a_{ij} = A^\text{free}_{ij} \quad (i \neq j)
  \label{eq:a-param}
\end{equation}
ensuring negative self-connections (default $-0.5$~Hz when the free
parameter is zero).

\subsubsection{Hemodynamic Model}

The Balloon-Windkessel model \citep{stephan2007} translates neural
activity into hemodynamic responses through four coupled ordinary
differential equations per region, implemented in log-space following
SPM12 (\texttt{spm\_fx\_fmri.m}):
\begin{align}
  \frac{ds}{dt} &= x - \kappa s - \gamma (f - 1)
    \label{eq:vasodilatory} \\
  \frac{d \ln f}{dt} &= \frac{s}{f}
    \label{eq:flow} \\
  \frac{d \ln v}{dt} &= \frac{f - v^{1/\alpha}}{\tau v}
    \label{eq:volume} \\
  \frac{d \ln q}{dt} &= \frac{f\, E(f, E_0) / E_0 - v^{1/\alpha} q / v}{\tau q}
    \label{eq:deoxyhemoglobin}
\end{align}
where $E(f, E_0) = 1 - (1 - E_0)^{1/f}$ is the oxygen extraction
fraction.

\begin{table}[h]
  \centering
  \caption{Hemodynamic parameter defaults (SPM12 code).}
  \label{tab:hemo-params}
  \begin{tabular}{llccl}
    \toprule
    Parameter & Symbol & Default & Unit & Description \\
    \midrule
    Signal decay       & $\kappa$ & 0.64 & s$^{-1}$ & Rate of vasodilatory signal decay \\
    Autoregulation     & $\gamma$ & 0.32 & s$^{-1}$ & Flow-dependent elimination \\
    Transit time       & $\tau$   & 2.00 & s        & Hemodynamic transit time \\
    Grubb's exponent   & $\alpha$ & 0.32 & --       & Vessel stiffness \\
    O$_2$ extraction   & $E_0$    & 0.40 & --       & Resting oxygen extraction fraction \\
    \bottomrule
  \end{tabular}
\end{table}

\subsubsection{BOLD Signal}

The BOLD percent signal change is computed from hemodynamic states via
the simplified Buxton observation equation \citep{stephan2007}:
\begin{equation}
  y = V_0 \left[ k_1 (1 - q) + k_2 \left(1 - \frac{q}{v}\right) + k_3 (1 - v) \right]
  \label{eq:bold}
\end{equation}
where $k_1 = 7 E_0$, $k_2 = 2$, $k_3 = 2 E_0 - 0.2$, and
$V_0 = 0.02$ is the resting venous blood volume fraction.

\subsubsection{Spectral DCM}

For resting-state fMRI, spectral DCM operates in the frequency domain
using cross-spectral density (CSD) as the data representation
\citep{friston2014}.

The transfer function is computed via eigendecomposition:
\begin{equation}
  H(\omega) = C_\text{out} (i\omega I - A)^{-1} C_\text{in}
    = \sum_k \frac{\mathbf{g}_k \mathbf{u}_k^T}{i 2\pi\omega - \lambda_k}
  \label{eq:transfer}
\end{equation}
where $\lambda_k$ are the eigenvalues of $A$ (stabilized by clamping
$\text{Re}(\lambda_k) \leq -1/32$) and $C_\text{in} = C_\text{out} = I_N$
in the standard spDCM convention.

The predicted CSD is:
\begin{equation}
  S(\omega) = H(\omega)\, G_u(\omega)\, H(\omega)^\dagger + G_n(\omega)
  \label{eq:predicted-csd}
\end{equation}
where the neuronal noise follows a $1/f^\alpha$ power law:
\begin{equation}
  G_{u,i}(\omega) = C \cdot \exp(a_{0,i}) \cdot \omega^{-\exp(a_{1,i})} \cdot 4
  \label{eq:neuronal-noise}
\end{equation}
with $C = 1/256$ (SPM scaling constant), and the observation noise
combines global and regional components:
\begin{equation}
  G_n(\omega) = G_\text{global}(\omega) \cdot \mathbf{1}\mathbf{1}^T
    + \text{diag}(G_\text{regional}(\omega))
  \label{eq:obs-noise}
\end{equation}

\subsubsection{Regression DCM}

Regression DCM \citep{frassle2017} reformulates DCM as a per-region
linear regression in the frequency domain:
\begin{equation}
  y_j = X_j \theta_j + \varepsilon_j
  \label{eq:rdcm-regression}
\end{equation}
where $y_j$ is the DFT of region $j$'s BOLD signal (real and imaginary
parts stacked), $X_j$ is the frequency-domain design matrix containing
HRF-convolved regressors from other regions and external inputs, and
$\theta_j$ contains the connectivity and input parameters for region $j$.

\subsection{Inference}

\subsubsection{Stochastic Variational Inference}

We replace the Variational Laplace algorithm used in SPM12
\citep{friston2007vfe} with stochastic variational inference (SVI;
\citealp{blei2017}) implemented in Pyro \citep{bingham2019}. The ELBO
objective is:
\begin{equation}
  \mathcal{L}(\phi)
    = \mathbb{E}_{q_\phi(\theta)}
      \left[ \log p(y, \theta \mid m) \right]
    - \mathbb{E}_{q_\phi(\theta)}
      \left[ \log q_\phi(\theta) \right]
  \label{eq:elbo}
\end{equation}
where the variational family is a mean-field Gaussian:
\begin{equation}
  q_\phi(\theta) = \prod_i \mathcal{N}(\theta_i;\, \mu_i, \sigma_i^2)
  \label{eq:mean-field}
\end{equation}
Optimization uses ClippedAdam with initial learning rate 0.01, gradient
clipping norm 10.0, and exponential learning rate decay to 1\% over the
training run. ODE-based models use fixed-step rk4 integration with
$\Delta t = 0.5$~s during SVI.

Priors: $A_\text{free} \sim \mathcal{N}(0, 1/64)$ (SPM12 convention),
$C \sim \mathcal{N}(0, 1)$.

\subsubsection{Amortized Inference}

For rapid posterior estimation across subjects, we provide amortized
inference using Neural Spline Flows \citep{papamakarios2021} conditioned
on summary statistics \citep{cranmer2020}:
\begin{equation}
  q_\phi(\theta \mid \mathbf{x})
    = f_\phi\!\left(\mathbf{z};\, \text{summary}(\mathbf{x})\right)
  \label{eq:amortized}
\end{equation}
Summary networks compress BOLD time series (1D-CNN) or CSD matrices
(MLP) into fixed-dimensional embeddings. Parameters are packed into a
standardized vector on $[-5, 5]$ matching the NSF spline domain.

\subsubsection{Analytic VB for rDCM}

Regression DCM admits closed-form posterior computation
\citep{frassle2017}:
\begin{align}
  \Sigma_j &= \left(\Lambda_0 + \langle\lambda_j\rangle X_j^H X_j\right)^{-1}
    \label{eq:rdcm-cov} \\
  \mu_j &= \Sigma_j
    \left(\Lambda_0 m_0 + \langle\lambda_j\rangle X_j^H y_j\right)
    \label{eq:rdcm-mean}
\end{align}
with noise precision updates $a_N = a_0 + N_y/2$ and
$b_N = b_0 + \frac{1}{2}(y^H y - \mu^T \Sigma^{-1} \mu + m_0^T \Lambda_0 m_0)$.

The free energy is computed per region as the sum of five terms
(log-likelihood, prior complexity, posterior entropy, noise prior, noise
posterior). The sparse variant \citep{frassle2018} extends this with ARD
binary indicators $z_{j,k} \sim \text{Bernoulli}(0.5)$.

\subsection{Bayesian Model Comparison}

For task and spectral DCM, the final ELBO serves as a lower bound on the
log model evidence \citep{friston2007vfe}. The model with the higher
ELBO is preferred. For rDCM, the analytic free energy $F = \sum_r F_r$
provides an exact bound within the variational approximation
\citep{frassle2017}. A difference $\Delta > 3$~nats (Bayes factor $> 20$)
is considered strong evidence.

\subsection{Implementation}

Pyro-DCM is implemented in Python 3.10+ using PyTorch 2.x for tensor
computation, Pyro 1.9+ for probabilistic programming, torchdiffeq for
ODE integration, and Zuko for normalizing flows. Numerical stability is
ensured via log-space hemodynamic states, eigenvalue stabilization
($\text{Re}(\lambda) \leq -1/32$), gradient clipping (norm 10.0),
NaN detection with early termination, and \texttt{torch.linalg.solve}
for matrix operations.

ODE integration uses adaptive dopri5 for simulation ($\Delta t = 0.01$~s)
and fixed-step rk4 for SVI ($\Delta t = 0.5$~s) to ensure predictable
runtime during optimization.

\subsection{Benchmarks}

Parameter recovery benchmarks follow a simulate--infer--compare protocol
with 3-region networks. Metrics include RMSE$(A)$, 90\% coverage,
Pearson correlation, ELBO/free energy, and wall time per subject.
Amortization-specific metrics include the RMSE ratio, amortization gap
($\text{ELBO}_\text{SVI} - \text{ELBO}_\text{amortized}$), and speed
ratio. See the benchmark report for detailed results.
